% numerals-ok: begin
\section{Three propositions and the seven tests between them}
\label{app:propositions}

The body reports seven negative results in seven places. One at a time they look like failures;
together they are one proposition tested seven times. This appendix states it, proves it, and
lines the tests up against it; two further propositions do the same for the reservation and for
the choice of its level.

\subsection*{Proposition 1 (the information channel)}

\begin{quote}
\itshape
Within the protocol's action set and information set, the data enter the objective only through
the duration law $F_d$ and the release sequence. Any change that leaves the posterior of the
source duration $d$ unchanged is, in expectation, the identity.
\end{quote}

\noindent\emph{Proof.}
(i) \emph{The realisation is deterministic.} Given a draw $(d_1,\dots,d_n)$ and a release
sequence, the service time $p_j=\alpha d_j+\beta$ is fixed, and start and completion times follow
from the policy's decisions and $s_{j+1}\ge s_j+\max(p_j,b_j)$; the weighted on-time rate is
therefore a deterministic function of $(d,\text{releases},\pi)$, with no other random source.
(ii) \emph{The data have one entrance.} A policy sees only $o_{1:t}$. Write $q_t(d)=P(d\mid o_{1:t})$
for the posterior. If an intervention leaves $q_t$ unchanged for every $t$ and every task family,
the policy sees the same information at every step, takes the same decisions, and reproduces the
same rate on every draw, hence the same expectation.
(iii) \emph{Sufficiency.} When the prefix is uninformative, $q_t\equiv F_d$, the optimal policy
depends only on the prior and on the release and due sequence. More features then have no
entrance: the channel is absent, not narrow.

\subsection*{The seven tests}

Table~\ref{tab:negatives} in Appendix~E is the audit trail. Each row is an intervention a
reviewer would expect to help; each changes the posterior of $d$ by less than the measurement
noise: the missing-frame channel (local fill, interpolation), reliability scores as sample
weights, self-supervised pre-training, a longer released prefix, a paid second observation, a
deeper sequence model, and the motion features of all 184 videos. The seventh row is the
sharpest, because it has a control: permuting the video columns reproduces the gain that the
video block appeared to buy, so the block was a wider design matrix, not information.

The upper bound is measured the same way. Giving the policy the true duration inside the same
rule is worth \slotPOracleGapSTwoPt{} points in S\slotCaseTwo{} and \slotPOracleGapSThreePt{} points
in S\slotCaseThree{} --- what perfect information about $d$ is worth here; every feature-side
intervention is bounded above by it, and none approaches it.

\subsection*{Proposition 2 (reservation neutrality)}

\begin{quote}
\itshape
Given the start time $s$, the on-time probability of a job is
$P(C\le \mathrm{due}\mid s)=F_d\bigl((\mathrm{due}-s-\beta)/\alpha\bigr)$, independent of the
reservation $b$. The reservation acts on the next job through $s'=s+\max(p,b)$ and on nothing else.
\end{quote}

\noindent\emph{Proof.} $C=s+p=s+\alpha d+\beta\le\mathrm{due}$ if and only if
$d\le(\mathrm{due}-s-\beta)/\alpha$; the right-hand side does not contain $b$, so the probability is
$F_d$ evaluated there.

The condition on $s$ cannot be dropped. If $b>p$ the next job starts later, its own
$(\mathrm{due}'-s'-\beta)/\alpha$ shrinks, and its on-time probability falls: the reservation
trades this job's capacity for the next job's time. In the most loaded case raising the reserve
from bottom to top costs about ten points of weighted on-time rate (Section~\ref{sec:p3res}),
while under the nominal mapping, where $b\le p$ throughout, all policies sit at the same rate
with reserved idle near zero.

\subsection*{Proposition 3 (reservation dominance)}

\begin{quote}
\itshape
Fix a policy whose choice of job does not depend on the reserve, and fix the duration
realisation. If $b'\le b$, then every start time under $b'$ is no later than under $b$, every
completion is no later, and every cost is no larger. Hence the smallest admissible reserve
weakly dominates every larger one; with the protocol floor $\underline{b}=0.2$,
$b^{\star}=\underline{b}$.
\end{quote}

\noindent\emph{Proof.} Order the executed jobs by start. The first start $s_1$ is fixed by the
release sequence; its completion $C_1=s_1+p_1$ contains no $b$. The next start is
$s_2=s_1+\max(p_1,b_1)$, which is non-decreasing in $b_1$; by induction every later start and
completion is non-decreasing in the whole reserve vector. A completion that is on time under $b$
therefore remains on time under $b'\le b$, so the weighted on-time set can only grow, while
$c=2+0.1b+0.2D$ can only fall and the feasibility test $s+b\le H$ can only become easier. Taking
expectations over the duration law preserves the order. If $b\le p$ for every job, then
$\max(p,b)=p$ and the bound is exactly flat; the plateau and the fall are the two branches of the
same statement.

\subsection*{What the three propositions say together}

The ordering they imply is the paper's one-sentence answer: the on-time rate is bought by the
rule, not by the data. Changing the rule moves the rate from \slotPImpSTwoAllBase{} to
\slotPImpSTwoAllCand{} in S\slotCaseTwo{} and from \slotPImpSThreeAllBase{} to
\slotPImpSThreeAllCand{} in S\slotCaseThree{}; knowing the true duration is worth
\slotPOracleGapSTwoPt{} and \slotPOracleGapSThreePt{} points; buying observation is worth no more
than zero; and adjusting the reservation is worth zero conditional on the start time, is worth
zero again while it stays below the service time, and turns negative once it exceeds it. That
order is the design principle the paper argues for; every acceptance bar behind it was written
down before the run (Table~\ref{tab:pass}).
% numerals-ok: end
